% =====================================================================
% CSCI 5360 Web Technologies -- Final Project Report (living document)
% Mathew Shereni -- Reliability of Serverless Web Applications (Blue Track)
%
% This file grows a little after every working day. Rules:
%   * Write facts while they are fresh: what was built, what was measured,
%     what changed from the midterm plan and why, and what was learned.
%   * Red \todo{...} text marks sections still to be written.
%   * Section numbering follows the midterm so each part can be traced back.
% =====================================================================
\documentclass[10pt]{article}
\usepackage[margin=0.75in]{geometry}
\usepackage{mathptmx}
\usepackage{microtype}
\usepackage{booktabs}
\usepackage{tabularx}
\usepackage{ragged2e}
\usepackage{graphicx}
\usepackage{enumitem}
\usepackage{xcolor}
\usepackage{amsmath}
\usepackage[font=small,skip=3pt]{caption}
\usepackage{titlesec}
\usepackage[hidelinks]{hyperref}

\setlist{nosep, topsep=2pt, leftmargin=1.5em}
\titlespacing*{\section}{0pt}{8pt plus 2pt minus 2pt}{3pt plus 1pt}
\titlespacing*{\subsection}{0pt}{5pt plus 2pt minus 1pt}{2pt plus 1pt}
\titleformat*{\section}{\large\bfseries}
\titleformat*{\subsection}{\normalsize\bfseries}
\renewcommand{\arraystretch}{0.92}

\newcommand{\todo}[1]{\textcolor{red}{\textbf{[TODO:} #1\textbf{]}}}

\begin{document}
\begin{center}
{\Large\bfseries Reliability of Serverless Web Applications --- Blue Track}\\[2pt]
{\large Final Project Report}\\[2pt]
Mathew Shereni \quad Banner ID 001460357\\
CSCI 5360 Web Technologies, Fall 2026 \quad December 10, 2026
\end{center}

% =====================================================================
\section{Assigned Project and Final Scope}
\todo{Objective as delivered; which of tasks T1--T5 were completed, each
with its completion criterion and evidence (commit, test, or figure).}

% =====================================================================
\section{Implementation}

\subsection{Development Environment and Safeguards}
Development began on October 30, 2026, on Ubuntu 26.04 with the tool
versions in Table~\ref{tab:tools}. Ubuntu 26.04 ships Python 3.14, while the
Lambda functions run on the Python 3.12 runtime fixed in the design. To keep
local tests and the cloud runtime on the same language version, the project
uses uv to install Python 3.12 next to the system interpreter and records the
version in a \texttt{.python-version} file; test dependencies are pinned in
\texttt{requirements-dev.txt}.

\begin{table}[h]
\centering
\footnotesize
\begin{tabular}{@{}ll@{}}
\toprule
\textbf{Tool} & \textbf{Version} \\
\midrule
Python (project environment, via uv) & 3.12.16 \\
AWS CLI & 2.31.35 \\
Terraform & 1.17.1 \\
k6 (local tests) & 2.3.0 \\
uv & 0.13.6 \\
\bottomrule
\end{tabular}
\caption{Development tool versions.}
\label{tab:tools}
\end{table}

All AWS work uses an IAM user with multi-factor authentication instead of the
root account, which is also protected by multi-factor authentication. The AWS
CLI uses a named profile (\texttt{b10}) fixed to us-east-2, and credentials are
stored outside the repository. A \$20 monthly cost budget sends email alerts;
it is configured to exclude credits, so that alerts reflect actual usage rather
than spending after the \$119 credit has been applied (Risk~1 in the midterm).
Because service quotas are set per region, the Lambda concurrency quota was
checked in us-east-2 itself with \texttt{aws lambda get-account-settings},
which confirmed 1{,}000 concurrent executions.

The public repository mirrors the module structure of the design:
\texttt{infra/} (M8), \texttt{api/} (M1--M3), \texttt{dep/} (M4),
\texttt{load/} (M5), \texttt{orchestrate/} (M6), \texttt{analysis/} (M7), and
\texttt{tests/}. The \texttt{.gitignore} file excludes Terraform state, local
experiment data, and credentials, while the Terraform lock file is kept under
version control so that the AWS provider version is pinned.

\todo{Following days: Terraform foundation, dependency function and fault
injector, Catalog API, reliability middleware, workload, orchestrator,
metrics.}

% =====================================================================
\section{Architecture as Built}
\todo{Final architecture and internal design figures; note any differences
from midterm Figures 1 and 2.}

% =====================================================================
\section{Methodology as Executed}
\todo{Research question, hypotheses, conditions, and repetitions as actually
run; deviations from the midterm design.}

% =====================================================================
\section{Testing}
\todo{Results of correctness tests CT1--CT7 (and T4/T5 tests if built).}

% =====================================================================
\section{Results}
\todo{SLO verdict matrix; latency and error tables; outage time series;
cache and concurrency figure; cost table.}

% =====================================================================
\section{Discussion}
\todo{Answers to RQ1--RQ3; status of H1--H3; competing explanations;
limitations.}

% =====================================================================
\section{Differences from the Midterm Plan}

\begin{table}[h]
\centering
\footnotesize
\begin{tabularx}{\linewidth}{@{}l>{\RaggedRight}X>{\RaggedRight\arraybackslash}X@{}}
\toprule
\textbf{Date} & \textbf{Change} & \textbf{Reason} \\
\midrule
Oct 30 & Milestone 1 (testbed, T1) moved from November 3 to November 8;
  a setup phase (October 30 to November 1) was added. &
  Implementation started on October 30, after the midterm submission, so the
  first milestone was rescheduled; later milestones keep their dates. \\
\bottomrule
\end{tabularx}
\caption{Changes from the midterm plan.}
\label{tab:changes}
\end{table}

% =====================================================================
\section{Lessons Learned}
\begin{itemize}
  \item Service quotas are regional: a quota increase must be checked in the
  region where the experiments run, not only where it was requested.
  \item Credits can hide spending: a budget that includes credits would show
  no cost while the credit is being used up, so the alert must exclude them.
  \item The local interpreter can differ from the cloud runtime (Python 3.14
  locally, 3.12 in Lambda); pinning the project version avoids tests that pass
  locally but fail in Lambda.
\end{itemize}

% =====================================================================
\section{AI Use Disclosure}
\begin{itemize}
  \item Oct 30--31: Claude (Anthropic) produced the implementation plan and
  explained and provided the setup commands; I typed and ran each command and
  checked its output before continuing.
\end{itemize}

% =====================================================================
\section*{References}
\todo{Carry over the verified references from the midterm and add any new
sources actually used.}

\end{document}
